\documentclass[11pt]{article}
% ---------------------------------------------------------------------------------------------
% Discussion paper: Numerical Intelligence Transfer.
% Every experimental number is read from numbers.tex, which src/export_latex.py writes from
% results/metrics.json (main.py runs it automatically at the end of a run). Missing values print
% as a red [TBD], so the paper compiles before and after the experiments.
% ---------------------------------------------------------------------------------------------
\usepackage[margin=1in]{geometry}
\usepackage[T1]{fontenc}
\usepackage[utf8]{inputenc}
\usepackage{lmodern}
\usepackage{microtype}
\usepackage{amsmath,amssymb}
\usepackage{booktabs}
\usepackage{array}
\usepackage{multirow}
\usepackage{graphicx}
\usepackage[table]{xcolor}
\usepackage{enumitem}
\usepackage[numbers,sort&compress]{natbib}
\usepackage{hyperref}
\usepackage{caption}
\usepackage{tcolorbox}
\setlength{\emergencystretch}{1.5em}
\hypersetup{colorlinks=true,linkcolor=blue!50!black,citecolor=blue!50!black,urlcolor=blue!50!black}

% ----------------------------------------------------------------- result macros
% Real results (numbers.tex) take precedence; illustrative dummies are used only if it is absent.
\newcommand{\dummyval}[1]{\textcolor{blue!70!black}{#1}}% set to {#1} to print dummies in black
\InputIfFileExists{numbers.tex}{}{\InputIfFileExists{dummy_numbers.tex}{}{}\InputIfFileExists{partial_real_numbers.tex}{}{}}% real values available so far override the dummies
\newcommand{\tbd}{\textcolor{red}{\textbf{[TBD]}}}
\makeatletter
\newcommand{\res}[2]{\@ifundefined{res@#1@#2}{\tbd}{\@nameuse{res@#1@#2}}}
\newcommand{\respm}[2]{\@ifundefined{res@#1@#2}{\tbd}{\@nameuse{res@#1@#2}%
  \@ifundefined{resstd@#1@#2}{}{\,{\scriptsize$\pm$\@nameuse{resstd@#1@#2}}}}}
\newcommand{\teach}[1]{\@ifundefined{teach@#1}{\tbd}{\@nameuse{teach@#1}}}
\newcommand{\envv}[1]{\@ifundefined{env@#1}{\tbd}{\@nameuse{env@#1}}}
\makeatother
\newcommand{\interp}[1]{\textcolor{red}{\textit{[Interpretation pending results: #1]}}}
\newcommand{\figorbox}[2]{\IfFileExists{figures/#1}{\includegraphics[width=#2]{figures/#1}}%
  {\fbox{\parbox[c][4cm][c]{\dimexpr#2-2\fboxsep-2\fboxrule\relax}{\centering\textcolor{red}{\textbf{[Figure pending]}}\\\scriptsize\texttt{\detokenize{#1}}}}}}

\newcommand{\Lq}{\mathcal{L}}
% ----------------------------------------------------------------- display names (rename here)
\usepackage{xspace}
\newcommand{\method}{\textup{\textsc{NumBridge}}\xspace}
\newcommand{\vBase}{\textup{\textsc{Base}}\xspace}                      % original all-MiniLM-L6-v2
\newcommand{\vParse}{\textup{\textsc{Base+Parse}}\xspace}                % original + symbolic parser channel
\newcommand{\vParseOnly}{\textup{\textsc{ParseOnly}}\xspace}             % parser alone (reference)
\newcommand{\vExact}{\textup{\textsc{NumBridge-Exact}}\xspace}           % exact (mathematical) supervision
\newcommand{\vMatched}{\textup{\textsc{NumBridge-Matched}}\xspace}       % exact labels on the reasoner's pair budget
\newcommand{\vRaw}{\textup{\textsc{NumBridge-Raw}}\xspace}               % raw LLM-reasoner labels
\newcommand{\vVerified}{\textup{\textsc{NumBridge-Verified}}\xspace}     % LLM-reasoner labels, verified
\newcommand{\vExactParse}{\textup{\textsc{NumBridge-Exact+Parse}}\xspace}% exact + parser channel
\newcommand{\vDual}{\textup{\textsc{NumBridge-Dual}}\xspace}             % learned numerical subspace, no parser
\newcommand{\vAbs}{\textup{\textsc{NumBridge-Abs}}\xspace}               % exact, absolute-distance target

\newcommand{\sem}{\mathrm{sem}}

% ----------------------------------------------------------------- config shorthands
\newcommand{\Ea}{exp1_original_minilm}
\newcommand{\Eb}{exp2_parser_plus_original}
\newcommand{\Ep}{ref_parser_only}
\newcommand{\Ec}{exp3_math_full}
\newcommand{\Ecb}{exp3b_math_budget}
\newcommand{\Eda}{exp4a_teacher_raw}
\newcommand{\Edb}{exp4b_teacher_verified}
\newcommand{\Ee}{exp5_hybrid_finetuned_plus_parser}
\newcommand{\Ef}{exp6_learned_projection}
\newcommand{\Eabs}{exp3_math_full_absdist}

\title{\textbf{NumBridge:\\Can Small LLMs Teach Sentence Encoders to Understand Numbers?}\\[4pt]
\large A Discussion Paper with Preliminary Results}
\author{Tariq Mahmood\\
\small AI \& Data Science Lab, DIMES\\Dept.\ of Computer Science, Modelling, Electronics and Systems Engineering\\
University of Calabria, Rende (CS), Italy}
\date{Draft --- \today}

\begin{document}
\maketitle

\ifdefined\ResultsAreDummy
\begin{tcolorbox}[colback=orange!10,colframe=orange!80!black,title=Illustrative dummy values]
Values shown in \dummyval{blue} are \textbf{placeholders for layout only, not experimental results}. Values in
black (reasoner accuracy, hardware) are real, from the run in progress. They are replaced automatically by the real values in \texttt{numbers.tex}. Remove this
build from circulation before submission.
\end{tcolorbox}
\else\ifdefined\ResultsAreMock
\begin{tcolorbox}[colback=red!8,colframe=red!70!black,title=Pipeline-check numbers]
The values in this build come from a pipeline validation run (randomly initialised tiny encoder, mock reasoner).
They verify that the code runs end to end and \textbf{are not findings} about all-MiniLM-L6-v2 or Qwen2.5.
\end{tcolorbox}
\else\ifdefined\ResultsAvailable\else
\begin{tcolorbox}[colback=blue!5,colframe=blue!50!black,title=Status of this draft]
The method, benchmark and experimental design are final. Experimental values marked \tbd{} and figures marked
\emph{[Figure pending]} are filled in automatically from \texttt{results/metrics.json} once the experiments
finish. Dataset statistics are final: the generator is deterministic.
\end{tcolorbox}
\fi\fi\fi

\begin{abstract}
Dense sentence encoders are the retrieval backbone of retrieval-augmented generation, semantic search and entity
resolution, yet they largely treat numbers as surface tokens. In such encoders \emph{10\,kg} and \emph{10{,}000\,g} may be far
apart, while \emph{10\,kg} and \emph{10\,m} look alike. Small instruction-tuned language models show some numerical
competence, but they are too expensive to run on every comparison during retrieval. This discussion paper asks whether
that competence can be \emph{transferred} into a lightweight encoder (all-MiniLM-L6-v2) by knowledge distillation,
and, crucially, how one would know that the encoder has learned numbers rather than benefiting from a symbolic parser.
We contribute (i) a formulation of what ``numerical intelligence'' can and cannot mean for a cosine-similarity space;
(ii) a distillation protocol in which a small LLM reasoner (Qwen2.5) is verified against deterministic ground truth,
paired with a budget-matched reasoner-free control that isolates what the reasoner actually adds; (iii) attribution
controls that separate embedding-internal learning from parser-driven gains, namely a parser-free numbers-in-words
slice and a learned numerical projection head; and (iv) a leakage-controlled synthetic benchmark of 5{,}010 statements
over six physical dimensions, with affine temperature handling. We argue that the most informative outcome is not
whether a single number improves. It is the pattern across controls. That pattern tells us whether the gains live in
the geometry, in the parser, or in the reasoner. We pre-declare hypotheses and decision rules, and we report
preliminary results under that protocol.
\end{abstract}

% =============================================================================================
\section{Introduction}
% =============================================================================================
Consider three statements: ``The package weighs 10\,kg'', ``The package weighs 10{,}000\,g'' and ``The package is
10\,m long''. The first two describe the same physical quantity; the third describes a different dimension altogether.
A retrieval system built on a general-purpose sentence encoder~\cite{reimers2019sbert} typically scores the third
statement as \emph{at least as similar} to the first as the second. The wording overlaps and the digits match, and the
encoder has no notion of unit conversion. Probing studies show that magnitude is partly decodable from token and
contextual embeddings but extrapolates poorly~\cite{naik2019numeracy,wallace2019numbers,zhang2020scales}; see
\citet{thawani2021survey} for a survey. In retrieval-augmented generation~\cite{lewis2020rag} and in entity
resolution over records with numeric attributes, such errors propagate silently. A wrong passage or a wrong match is
retrieved, and the downstream model reasons correctly over incorrect evidence.

Small instruction-tuned language models, such as the Qwen2.5 family~\cite{qwen2025qwen25}, can often answer ``is
10\,kg the same as 10{,}000\,g?'' correctly. Running such a model inside every similarity computation is impractical,
however. This motivates \textbf{Numerical Intelligence Transfer}: distilling the numerical behaviour of a small LLM
reasoner into a compact encoder~\cite{hinton2015distilling,reimers2020multilingual}, so that the numerical
competence is paid for once, at training time. We call the resulting framework \method: it bridges the numerical competence of an LLM \emph{reasoner} and the efficiency of a compact retrieval \emph{encoder}.

The idea is simple; evaluating it honestly is not. We identify three confounds that can make a transfer method look
successful when it is not:
\begin{enumerate}[leftmargin=*,itemsep=2pt]
\item \textbf{Cosine is not numeracy.} A higher cosine for equivalent pairs is meaningless if every pair's cosine
rises. Similarity has to be \emph{discriminative} and \emph{graded}, not merely high.
\item \textbf{The parser confound.} Any pipeline that extracts numbers symbolically will excel on text it can parse.
Gains can come from the parser while the embedding learns nothing.
\item \textbf{The verification paradox.} If reasoner labels are filtered against mathematical ground truth, which is
good practice, then on every determinable relation the surviving labels \emph{equal} that ground truth. What, then,
does the reasoner contribute beyond direct mathematical supervision?
\end{enumerate}

\paragraph{Contributions.} (1) A formulation of numerical desiderata for embedding spaces, including what symmetric
cosine geometry cannot express (\S\ref{sec:formulation}). (2) A verified-reasoner distillation protocol with five
objectives and three representation choices (\S\ref{sec:method}). (3) An experimental design whose controls target
the three confounds: a budget-matched reasoner-free control, a parser-free stress slice, a learned projection head, and
quantity-level leakage control (\S\ref{sec:design}). (4) Pre-declared hypotheses and decision rules, with preliminary
results (\S\ref{sec:results}). (5) A discussion of open problems and of where LLM reasoners can \emph{genuinely} add
information (\S\ref{sec:discussion}).

% =============================================================================================
\section{Background and Related Work}
% =============================================================================================
\paragraph{Numeracy in learned representations.} \citet{naik2019numeracy} and \citet{wallace2019numbers} probe word
and contextual embeddings for magnitude and ordering. They find partial but brittle numeracy, which degrades outside
the training range. \citet{zhang2020scales} show that embeddings encode physical scales only coarsely, and
NumerSense~\cite{lin2020numersense} probes numerical commonsense in masked language models. \citet{thawani2021survey}
review number representations and argue for numeracy-aware design.

\paragraph{Numeracy-preserving representations.} Explicit numeric encodings, such as DICE~\cite{sundararaman2020dice},
build numerical distance into the embedding by construction. Language-modelling work improves number
prediction~\cite{spithourakis2018numeracy}, and \citet{geva2020injecting} inject numerical reasoning skills through
synthetic data. These approaches target token or generative representations. We target \emph{sentence-level
retrieval embeddings}, where numbers must coexist with context and meaning.

\paragraph{Sentence embeddings and distillation.} Sentence-BERT~\cite{reimers2019sbert} and contrastive objectives such
as SimCSE~\cite{gao2021simcse} produce the encoders deployed today. MiniLM~\cite{wang2020minilm} is itself a distilled
model. \citet{reimers2020multilingual} distil an embedding space into a encoder while preserving it, which is the
template for our semantic-preservation term. Benchmarks such as MTEB~\cite{muennighoff2023mteb} and STS~\cite{cer2017sts}
measure general quality but contain few numerically discriminative cases.

\paragraph{Preserving pretrained behaviour.} Fine-tuning risks catastrophic forgetting. Parameter-space
regularisers such as EWC~\cite{kirkpatrick2017ewc} and L2-SP~\cite{li2018l2sp} constrain weights. We instead constrain
\emph{outputs}, distilling from the frozen pretrained encoder on general text.

\paragraph{Probing methodology and order.} Linear probes measure decodability, not use, and high probe accuracy can
reflect probe capacity~\cite{hewitt2019control}. We therefore report geometric (cosine) and probe metrics separately.
Order is antisymmetric, and order embeddings~\cite{vendrov2016order} encode it with asymmetric scores, a point that
matters because cosine is symmetric (\S\ref{sec:formulation}).

\paragraph{Positioning.} Prior work either \emph{probes} numeracy or \emph{adds} explicit numeric channels. We study
\emph{transfer} from an LLM reasoner into a retrieval encoder. The distinctive element is the experimental controls,
which attribute gains to the reasoner, the parser or the learned geometry.

% =============================================================================================
\section{What Can ``Numerical Intelligence'' Mean for an Embedding Space?}
\label{sec:formulation}
% =============================================================================================
Let each statement $x$ express a quantity $\kappa(x) = (d_x, v_x)$, where $d_x$ is its physical dimension (mass, length,
time, temperature, volume, proportion) and $v_x>0$ is its value in a canonical unit (g, m, s, K, L, dimensionless).
Temperature is \emph{affine}: $v = T_{\mathrm C} + 273.15$ and $v = (T_{\mathrm F}-32)\cdot\tfrac59 + 273.15$. Ratios of
Celsius or Fahrenheit readings are physically meaningless, so every ratio and log distance is computed on the
absolute scale. Let $e(x)\in\mathbb{S}^{383}$ be the normalised embedding and $s(x,y)=e(x)^\top e(y)$ its cosine.

We ask an embedding to satisfy five desiderata:
\begin{description}[leftmargin=*,itemsep=2pt]
\item[D1 Equivalence.] If $d_x=d_y$ and $v_x=v_y$, as with 10\,kg and 10{,}000\,g, or 10 and 10.00, then $s(x,y)$ should
exceed $s(x,z)$ for non-equivalent $z$ of the same dimension. Note that 10.01 is \emph{not} equivalent to 10.
\item[D2 Graded distance.] For same-dimension pairs, $1-s$ should increase monotonically with $|\log_{10}(v_x/v_y)|$. We
also examine relative absolute distance $|v_x-v_y|/\max(v_x,v_y)$. The two disagree, and which one is ``right'' depends
on the application.
\item[D3 Order.] The ordering $v_x<v_y$ should be recoverable from the representation.
\item[D4 Dimensional separation.] If $d_x\neq d_y$, then $s(x,y)$ should be low. 10\,kg and 10\,m are not comparable, so no
distance between them is meaningful.
\item[D5 Context and semantic preservation.] ``The patient's temperature is 37\,\textdegree C'' and ``the machine operates
at 37\,\textdegree C'' are numerically equal but contextually different. General semantic similarity should not
degrade.
\end{description}

\paragraph{Three tensions.} First, \textbf{D3 cannot be expressed by cosine}. Cosine is symmetric, so $s(x,y)$ says
nothing about which value is larger. Order can only be read off a direction or a probe; in cosine geometry it shows
up at best as \emph{betweenness}, where for $v_a<v_b<v_c$ we have $s(a,c)<\min(s(a,b),s(b,c))$. Second, \textbf{D1
competes with D5}. Two statements of the same quantity in different contexts must be ``equal'' numerically and
``different'' semantically, so a single scalar similarity has to compromise. Third, \textbf{D2 competes with D1 at
the top of the scale}. If near-equal pairs are pushed towards cosine~1, exact equivalence becomes harder to separate
from closeness. We therefore cap graded targets below a gap, $c_{\max}<1$.

These tensions suggest that a single cosine space may be the wrong object. One alternative is a \emph{separate
numerical subspace} learned beside the semantic one; another is a \emph{hybrid similarity} that adds a parser channel.
We test both.

% =============================================================================================
\section{Method}
\label{sec:method}
% =============================================================================================
\subsection{Verified reasoner supervision}
A small LLM reasoner, Qwen2.5-1.5B-Instruct~\cite{qwen2025qwen25} (Qwen2.5-0.5B-Instruct on limited hardware), answers
four questions about statement pairs and one about single statements (Appendix~\ref{app:prompts}). These are: same
dimension (Yes/No), equivalence after conversion (Yes/No), which is larger (A/B/Equal), the ratio $A/B$ after
conversion (a free-form number; asked for a random half of the same-dimension pairs), and whether a value is plausible in its context (Yes/No). Categorical questions are
answered in a single forward pass by reading the next-token distribution over the option tokens. This yields a hard
label and a soft confidence $p$, which is the distillation signal. Only the ratio question uses greedy decoding.
\textbf{Every determinable answer is checked against Python ground truth}, and reasoner accuracy is reported per
question. LLM output is never used as ground truth.

\subsection{Encoder objectives}
The encoder is all-MiniLM-L6-v2. It is trained with
\begin{equation}
\Lq = \lambda_1\Lq_{\mathrm{eq}} + \lambda_2\Lq_{\mathrm{ord}} + \lambda_3\Lq_{\mathrm{dist}} + \lambda_4\Lq_{\mathrm{ctx}} + \lambda_5\Lq_{\sem}.
\end{equation}
\textbf{Equivalence} uses a confidence-weighted symmetric InfoNCE~\cite{oord2018cpc} over same-context equivalent pairs
$(a_i,b_i)$ with weights $w_i$ (reasoner confidence, or 1 under direct supervision):
\begin{equation}
\Lq_{\mathrm{eq}} = -\frac{1}{2\sum_i w_i}\sum_i w_i\Big[\log\operatorname{softmax}_j\big(s(a_i,b_j)/\tau\big)_i + \log\operatorname{softmax}_j\big(s(b_i,a_j)/\tau\big)_i\Big].
\end{equation}
\textbf{Ordering} combines a margin ranking loss on a learned direction $u$, with $r(x)=u^\top e(x)$, and a betweenness
loss on ordered triplets $v_{\mathrm{lo}}<v_{\mathrm{mid}}<v_{\mathrm{hi}}$:
\begin{equation}
\Lq_{\mathrm{ord}} = \overline{[m_r-(r(x_{\mathrm{hi}})-r(x_{\mathrm{lo}}))]_+} + \overline{[s_{\mathrm{lo,hi}}-s_{\mathrm{lo,mid}}+m_t]_+ + [s_{\mathrm{lo,hi}}-s_{\mathrm{mid,hi}}+m_t]_+}.
\end{equation}
\textbf{Distance} regresses cosine onto a magnitude-similarity target for same-dimension, non-equivalent pairs:
\begin{equation}
\Lq_{\mathrm{dist}} = \overline{\big(s(a,b)-t(\Delta)\big)^2},\qquad t(\Delta)=c_{\min}+(c_{\max}-c_{\min})\,e^{-\Delta/\sigma},
\end{equation}
with $\Delta=|\log_{10}(v_a/v_b)|$ (logarithmic mode) or $\Delta=|v_a-v_b|/\max(v_a,v_b)$ (absolute mode). With the
reasoner, $\Delta$ is computed from the reasoner's ratio. \textbf{Context} places upper bounds on similarity:
$\Lq_{\mathrm{ctx}}=\overline{[s(a,b)-m_k]_+}$, where $m_k$ depends on whether the pair has different dimensions, is the
same quantity in a different context, or is a plausible versus contradictory statement. \textbf{Semantic preservation}
distils from the frozen pretrained encoder $f_0$ on unlabeled general sentences $g$:
$\Lq_{\sem}=\overline{1-e(g)^\top f_0(g)}$~\cite{reimers2020multilingual}.

\subsection{Representation choices}
\begin{itemize}[leftmargin=*,itemsep=2pt]
\item \textbf{Single space:} all objectives act on $e(x)$, and similarity is $s$.
\item \textbf{Learned numerical projection:} a linear head $P$ ($384\to128$) receives the numerical objectives on
$z(x)=\mathrm{norm}(Pe(x))$, while $\Lq_\sem$ acts on $e(x)$. Similarity is $\alpha s + (1-\alpha) z^\top z'$, with
$\alpha$ tuned on validation. No parser is involved.
\item \textbf{Parser hybrid:} a regex parser extracts $(d,v)$, and similarity is
$\alpha s + (1-\alpha)s_{\mathrm{num}}$. Here $s_{\mathrm{num}}=1$ for exact equivalence,
$c_{\max}e^{-\Delta/\sigma}$ for the same dimension, and $0$ across dimensions; it falls back to $s$ when parsing fails.
\end{itemize}

% =============================================================================================
\section{Benchmark and Experimental Design}
\label{sec:design}
% =============================================================================================
\subsection{A leakage-controlled synthetic benchmark}
The generator is deterministic (seed 42). It produces \textbf{5{,}010 statements} in seven categories
(Table~\ref{tab:data}) over six dimensions: length 1{,}184, volume 994, mass 897, temperature 885, time 843,
proportion 207. Formats include decimals, negative temperatures, thousands separators, scientific notation
(\texttt{2.5e6}) and abbreviated or spelled-out units across metric and imperial systems.

\textbf{Leakage control.} Every statement is keyed by its quantity $(d, v)$. All statements of the same quantity, in
any unit, precision, context or category, share an equivalence group and \emph{the same split}. Multi-quantity sets,
such as ordering triplets, are placed in one split whole. Seven held-out templates appear only in test (76 test
statements), giving an unseen-phrasing slice. The split is 3{,}541 / 761 / 708 (70.7\,/\,15.2\,/\,14.1\,\%), and zero
equivalence groups cross splits. Pairs and triplets are built within each split. The test set has 4{,}135 pairs
(351 same-context equivalent, 327 same quantity in a different context, 1{,}134 close ($\le2\times$), 597 mid, 1{,}047
distant ($\ge10\times$), 642 cross-dimension, 37 plausibility) and 2{,}600 ordering triplets.

\textbf{Parser-free stress slice.} 379 test statements are re-expressed with spelled-out integers (``thirty-five degrees
Celsius''). The parser recovers \emph{none} of them, so any gain on this slice must come from the embedding.

\begin{table}[t]
\centering\small
\caption{Benchmark categories. Ground truth is computed in Python for every relation.}
\label{tab:data}
\resizebox{\linewidth}{!}{%
\begin{tabular}{@{}llrl@{}}
\toprule
Cat. & Content & $n$ & Example \\
\midrule
A & Magnitude & 900 & \textit{The cable measures 0.75 miles.} \\
B & Unit equivalence sets & 1{,}202 & \textit{The swimming pool contains 15 L of water.} \\
C & Ordering triplets & 601 & \textit{The lab sample has a mass of 800,000 milligrams.} \\
D & Same quantity, many contexts & 602 & \textit{The insect is 25 feet long.} \\
E & Precision (10 / 10.0 / 10.00 / 10.01 / 10.1) & 500 & \textit{The bottle holds 12.0 centilitres.} \\
F & Ratios, percentages, differences & 604 & \textit{Half of 10 miles is 5 miles.} \\
G & Plausible vs contradictory & 601 & \textit{The child's syrup dose is 0.006 mL.} \\
\bottomrule
\end{tabular}}
\end{table}

\subsection{Configurations}
Table~\ref{tab:configs} lists the configurations. Three are controls designed for the confounds in the introduction.
\vMatched trains on \emph{exactly} the reasoner's pairs but with ground-truth labels, so \vVerified$-$\vMatched isolates what
the reasoner adds beyond its pair selection and confidence weights (the verification paradox). \vDual learns a
numerical subspace without any parser. The \textbf{words slice} disables the parser at test time (the parser confound). To keep this comparison fair, \vMatched mirrors the reasoner's supervision exactly: it receives distance targets only on the pairs for which the reasoner was asked a ratio (half of the same-dimension pairs), and neither it nor the reasoner-based models use the triplet betweenness term, which only \vExact receives. The same-/different-context labels used to form positives come from the benchmark metadata for every model, not from the reasoner.

\begin{table}[t]
\centering\small
\caption{Experimental configurations. All share the same test data; thresholds and $\alpha$ are tuned on validation.}
\label{tab:configs}
\begin{tabular}{@{}l>{\raggedright\arraybackslash}p{10.5cm}@{}}
\toprule
\vBase & Original all-MiniLM-L6-v2, cosine. \\
\vParse & Original embeddings + parser channel (hybrid similarity; probes see $[e;\text{parser features}]$). \\
\vParseOnly & Parser only: an upper reference on parser-readable text. \\
\vExact & Fine-tuned with direct mathematical supervision on all training pairs and triplets. \\
\vMatched & Direct supervision restricted to the reasoner's pair budget: the \textbf{reasoner-free control}. \\
\vRaw & Same pairs, raw reasoner labels, confidence-weighted. \\
\vVerified & Same pairs, reasoner labels kept only where verified correct. \\
\vExactParse & \vExact embeddings + parser channel. \\
\vDual & Learned numerical projection head (no parser). \\
\vAbs & \vExact with the absolute-distance target (seed 42 only). \\
Ablation & $\Lq_{\mathrm{sem}}$ only, then $+\Lq_{\mathrm{eq}}$, $+\Lq_{\mathrm{ord}}$, $+\Lq_{\mathrm{dist}}$, $+\Lq_{\mathrm{ctx}}$ (= \vExact). \\
\bottomrule
\end{tabular}
\end{table}

\subsection{Metrics}
\textbf{Equivalence}: ROC-AUC and balanced accuracy separating same-context equivalent pairs from \emph{hard}
negatives (same dimension, $\le10\times$ apart), with subsets for cross-unit conversion and precision.
\textbf{Ordering}: a per-dimension ridge probe to $\log_{10}v$, reporting pairwise order accuracy and Spearman $\rho$,
and betweenness accuracy under each model's similarity $s$. \textbf{Distance}: Spearman correlation between $1-s$ and $\Delta$ on
same-template pairs, which controls for context. \textbf{Retrieval}: every test statement queries all others, with
relevant items being the same equivalence group; we report R@$k$ and MRR. \textbf{Context}: dimension discrimination,
$\Pr[s(a,\text{equiv.})>s(a,\text{cross-dim.})]$; similarity of the same quantity in different contexts; linear
context and plausibility probes. \textbf{Semantics}: STS-B~\cite{cer2017sts} test Spearman.

\subsection{Pre-declared hypotheses and decision rule}
A difference is \emph{supported} if it exceeds $\max(0.01, 2\hat\sigma)$, where $\hat\sigma$ is the standard deviation
across seeds (seeds 42--44 for the main configurations). It is \emph{contradicted} if it is a decrease of that size,
and \emph{inconclusive} otherwise. We declare, before seeing results:
\begin{description}[leftmargin=*,itemsep=1pt]
\item[H1] The reasoner is reliable on dimension identification but less reliable on cross-unit equivalence and ratios.
\item[H2] The original encoder fails D1 and D2. Its cosine is high for all pair types, so equivalence AUC on hard
negatives is near chance, and its distance correlation is weak.
\item[H3] Direct supervision (\vExact) improves D1, D2 and D4 over \vBase on digit-form text.
\item[H4] Verified distillation (\vVerified) does not exceed budget-matched direct supervision (\vMatched), because of the
verification paradox; raw labels (\vRaw) do worse.
\item[H5] Parser hybrids (\vParse, \vExactParse) dominate on digit-form text but gain nothing on the words slice.
\item[H6] Semantic preservation keeps the STS-B drop within 0.02.
\end{description}

% =============================================================================================
\section{Preliminary Results}
\label{sec:results}
% =============================================================================================
\paragraph{Setup.} Hardware: \envv{hardware}; PyTorch \envv{torch}; total run time \envv{minutes} minutes. Reasoner:
\teach{teacher_model}. Means $\pm$ standard deviation over seeds are shown where more than one seed was run.

\subsection{Is the reasoner reliable? (H1)}
\begin{table}[h]
\centering\small
\caption{Reasoner accuracy against Python ground truth on \teach{n_pairs} training pairs.}
\label{tab:reasoner}
\begin{tabular}{@{}lr@{\hspace{2em}}lr@{}}
\toprule
Same dimension & \teach{same_dimension_acc} & Ordering & \teach{ordering_acc} \\
Equivalence (balanced) & \teach{equivalence_balanced_acc} & Ratio within 1\,\% & \teach{ratio_acc_within_1pct} \\
Equivalence AUC (soft) & \teach{equivalence_auc} & Ratio parse rate & \teach{ratio_parse_rate} \\
Cross-unit equiv.\ (recall) & \teach{unit_conversion_acc_cross_unit_equiv} & Plausibility (balanced) & \teach{compatibility_balanced_acc} \\
Majority baseline (equiv.) & \teach{majority_baseline_equivalence} & Fully verified pairs & \teach{fully_verified_fraction} \\
\bottomrule
\end{tabular}
\end{table}
The reasoner is far from reliable. Its equivalence judgements carry some signal (soft AUC \teach{equivalence_auc}),
but raw accuracy (\teach{equivalence_acc}) barely exceeds the majority baseline (\teach{majority_baseline_equivalence}) and
balanced accuracy is only \teach{equivalence_balanced_acc}. It recognises only \teach{unit_conversion_acc_cross_unit_equiv} of
cross-unit equivalences and places only \teach{ratio_acc_within_1pct} of ratios within 1\,\%, although every ratio answer
parses as a number. Its dimension and order answers collapse onto a single option: it answers ``Yes'' to the
same-dimension question for only 3\,\% of pairs, although 84\,\% share a dimension (accuracy \teach{same_dimension_acc}), and it
places 89\,\% of its order probability on ``Equal'' (accuracy \teach{ordering_acc}, close to the share of equal pairs).
Only \teach{fully_verified_fraction} of the pairs survive verification. H1 is therefore only \emph{partly} supported: the
reasoner is weak on conversion and ratios, as predicted, but it is \emph{not} reliable on dimension identification.
\textbf{This is a property of the model, not of the scoring.} A diagnostic confirms that 95--100\,\% of the
next-token probability falls on the option tokens, that free-text answers agree with the scored ones (``No'' for
10\,kg vs 10{,}000\,g; ``Equal'' for 2\,m vs 300\,cm), and that batching changes scores only at the third decimal.
Rewording the questions on 150 training pairs shifts \emph{which} option the model prefers (order answers move from
``Equal'' to mostly ``B'') but not its accuracy: 0.18 to 0.22 on dimension (majority baseline 0.86) and 0.31 to 0.33
on order (majority 0.39). A 1.5B reasoner therefore supplies little usable supervision on determinable relations,
which bears directly on H4.

\subsection{What the original encoder sees (H2)}
\begin{table}[h]
\centering\small
\caption{Mean cosine similarity by test pair type: original encoder (\vBase) vs direct fine-tuning (\vExact, seed 42).}
\label{tab:cos}
\begin{tabular}{@{}lccc@{}}
\toprule
Pair type & $n$ & \vBase & \vExact \\
\midrule
Equivalent, same context & \res{\Ea}{cosn.equivalent} & \res{\Ea}{cos.equivalent} & \res{\Ec}{cos.equivalent} \\
Same quantity, different context & \res{\Ea}{cosn.equivalent_diff_context} & \res{\Ea}{cos.equivalent_diff_context} & \res{\Ec}{cos.equivalent_diff_context} \\
Close ($\le 2\times$) & \res{\Ea}{cosn.close} & \res{\Ea}{cos.close} & \res{\Ec}{cos.close} \\
Mid & \res{\Ea}{cosn.mid} & \res{\Ea}{cos.mid} & \res{\Ec}{cos.mid} \\
Distant ($\ge 10\times$) & \res{\Ea}{cosn.distant} & \res{\Ea}{cos.distant} & \res{\Ec}{cos.distant} \\
Different dimension & \res{\Ea}{cosn.cross_dimension} & \res{\Ea}{cos.cross_dimension} & \res{\Ec}{cos.cross_dimension} \\
Plausible vs contradictory & \res{\Ea}{cosn.plausibility} & \res{\Ea}{cos.plausibility} & \res{\Ec}{cos.plausibility} \\
\bottomrule
\end{tabular}
\end{table}
\interp{is the \vBase column flat (all pair types similar)? If so, high cosine on equivalent pairs carries no numerical
information, and confound~1 holds.}

\begin{figure}[h]
\centering\figorbox{01_cosine_baseline_vs_finetuned.png}{0.95\linewidth}
\caption{Cosine similarity by numerical relation for the original, directly fine-tuned and reasoner-distilled encoders.}
\label{fig:cos}
\end{figure}

\subsection{Numerical equivalence, ordering and distance (H3, H4)}
\begin{table}[h]
\centering\small\setlength{\tabcolsep}{4pt}
\caption{Numerical metrics on the test split (digit-form text). AUC: equivalent vs hard negatives. Order: probe
pairwise accuracy. Betw.: betweenness under the model's similarity. $\rho_{\log}$: Spearman of $1-s$ vs $|\Delta\log_{10}|$ (same template).}
\label{tab:num}
\resizebox{\linewidth}{!}{%
\begin{tabular}{@{}lccccccc@{}}
\toprule
Model & Equiv.\ AUC & Unit conv. & Precision & Order (probe) & Betw. & $\rho_{\log}$ & MRR \\
\midrule
\vBase & \respm{\Ea}{equivalence.equivalence_auc_hard} & \respm{\Ea}{equivalence.unit_conversion_bal_acc} & \respm{\Ea}{equivalence.precision_bal_acc} & \respm{\Ea}{ordering.probe_pairwise_order_acc} & \respm{\Ea}{ordering.triplet_betweenness_acc} & \respm{\Ea}{distance.same_template_log_distance_spearman} & \respm{\Ea}{retrieval.mrr} \\
\vParse & \respm{\Eb}{equivalence.equivalence_auc_hard} & \respm{\Eb}{equivalence.unit_conversion_bal_acc} & \respm{\Eb}{equivalence.precision_bal_acc} & \respm{\Eb}{ordering.probe_pairwise_order_acc} & \respm{\Eb}{ordering.triplet_betweenness_acc} & \respm{\Eb}{distance.same_template_log_distance_spearman} & \respm{\Eb}{retrieval.mrr} \\
\vParseOnly & \respm{\Ep}{equivalence.equivalence_auc_hard} & \respm{\Ep}{equivalence.unit_conversion_bal_acc} & \respm{\Ep}{equivalence.precision_bal_acc} & \respm{\Ep}{ordering.probe_pairwise_order_acc} & \respm{\Ep}{ordering.triplet_betweenness_acc} & \respm{\Ep}{distance.same_template_log_distance_spearman} & \respm{\Ep}{retrieval.mrr} \\
\midrule
\vExact & \respm{\Ec}{equivalence.equivalence_auc_hard} & \respm{\Ec}{equivalence.unit_conversion_bal_acc} & \respm{\Ec}{equivalence.precision_bal_acc} & \respm{\Ec}{ordering.probe_pairwise_order_acc} & \respm{\Ec}{ordering.triplet_betweenness_acc} & \respm{\Ec}{distance.same_template_log_distance_spearman} & \respm{\Ec}{retrieval.mrr} \\
\vMatched & \respm{\Ecb}{equivalence.equivalence_auc_hard} & \respm{\Ecb}{equivalence.unit_conversion_bal_acc} & \respm{\Ecb}{equivalence.precision_bal_acc} & \respm{\Ecb}{ordering.probe_pairwise_order_acc} & \respm{\Ecb}{ordering.triplet_betweenness_acc} & \respm{\Ecb}{distance.same_template_log_distance_spearman} & \respm{\Ecb}{retrieval.mrr} \\
\vRaw & \respm{\Eda}{equivalence.equivalence_auc_hard} & \respm{\Eda}{equivalence.unit_conversion_bal_acc} & \respm{\Eda}{equivalence.precision_bal_acc} & \respm{\Eda}{ordering.probe_pairwise_order_acc} & \respm{\Eda}{ordering.triplet_betweenness_acc} & \respm{\Eda}{distance.same_template_log_distance_spearman} & \respm{\Eda}{retrieval.mrr} \\
\vVerified & \respm{\Edb}{equivalence.equivalence_auc_hard} & \respm{\Edb}{equivalence.unit_conversion_bal_acc} & \respm{\Edb}{equivalence.precision_bal_acc} & \respm{\Edb}{ordering.probe_pairwise_order_acc} & \respm{\Edb}{ordering.triplet_betweenness_acc} & \respm{\Edb}{distance.same_template_log_distance_spearman} & \respm{\Edb}{retrieval.mrr} \\
\vExactParse & \respm{\Ee}{equivalence.equivalence_auc_hard} & \respm{\Ee}{equivalence.unit_conversion_bal_acc} & \respm{\Ee}{equivalence.precision_bal_acc} & \respm{\Ee}{ordering.probe_pairwise_order_acc} & \respm{\Ee}{ordering.triplet_betweenness_acc} & \respm{\Ee}{distance.same_template_log_distance_spearman} & \respm{\Ee}{retrieval.mrr} \\
\vDual & \respm{\Ef}{equivalence.equivalence_auc_hard} & \respm{\Ef}{equivalence.unit_conversion_bal_acc} & \respm{\Ef}{equivalence.precision_bal_acc} & \respm{\Ef}{ordering.probe_pairwise_order_acc} & \respm{\Ef}{ordering.triplet_betweenness_acc} & \respm{\Ef}{distance.same_template_log_distance_spearman} & \respm{\Ef}{retrieval.mrr} \\
\vAbs & \respm{\Eabs}{equivalence.equivalence_auc_hard} & \respm{\Eabs}{equivalence.unit_conversion_bal_acc} & \respm{\Eabs}{equivalence.precision_bal_acc} & \respm{\Eabs}{ordering.probe_pairwise_order_acc} & \respm{\Eabs}{ordering.triplet_betweenness_acc} & \respm{\Eabs}{distance.same_template_log_distance_spearman} & \respm{\Eabs}{retrieval.mrr} \\
\bottomrule
\end{tabular}}
\end{table}
\interp{H3: \vExact vs \vBase per column, applying the decision rule. H4: \vVerified vs \vMatched, and \vRaw vs \vMatched. Note whether
ordering moves in the probe but not in betweenness, as \S\ref{sec:formulation} predicts.}

\begin{figure}[h]
\centering
\figorbox{02_numerical_vs_embedding_distance.png}{0.9\linewidth}
\caption{Numerical log distance vs embedding distance on same-template pairs.}
\label{fig:dist}
\end{figure}

\subsection{Context and semantics (H6)}
\begin{table}[h]
\centering\small
\caption{Context separation and semantic preservation.}
\label{tab:ctx}
\resizebox{\linewidth}{!}{%
\begin{tabular}{@{}lccccc@{}}
\toprule
Model & Dim.\ discrim. & Sim.\ same qty / diff.\ ctx & Context probe & Plausibility probe & STS-B $\rho$ \\
\midrule
\vBase & \respm{\Ea}{context.dimension_discrimination_acc} & \respm{\Ea}{context.same_quantity_diff_context_mean_sim} & \respm{\Ea}{context.context_probe_acc} & \respm{\Ea}{context.plausibility_probe_acc} & \respm{\Ea}{semantic.stsb_spearman} \\
\vExact & \respm{\Ec}{context.dimension_discrimination_acc} & \respm{\Ec}{context.same_quantity_diff_context_mean_sim} & \respm{\Ec}{context.context_probe_acc} & \respm{\Ec}{context.plausibility_probe_acc} & \respm{\Ec}{semantic.stsb_spearman} \\
\vVerified & \respm{\Edb}{context.dimension_discrimination_acc} & \respm{\Edb}{context.same_quantity_diff_context_mean_sim} & \respm{\Edb}{context.context_probe_acc} & \respm{\Edb}{context.plausibility_probe_acc} & \respm{\Edb}{semantic.stsb_spearman} \\
\vDual & \respm{\Ef}{context.dimension_discrimination_acc} & \respm{\Ef}{context.same_quantity_diff_context_mean_sim} & \respm{\Ef}{context.context_probe_acc} & \respm{\Ef}{context.plausibility_probe_acc} & \respm{\Ef}{semantic.stsb_spearman} \\
\bottomrule
\end{tabular}}
\end{table}
\interp{STS-B change for \vExact and \vDual relative to \vBase (H6); does the context loss keep same-quantity/different-context
similarity below equivalence?}

\subsection{Is it the embedding or the parser? (H5)}
\begin{table}[h]
\centering\small
\caption{Parser-free slice: numbers spelled out in words. The parser reads none of these statements, so hybrids fall
back to their embedding component.}
\label{tab:words}
\resizebox{\linewidth}{!}{%
\begin{tabular}{@{}lcccc@{}}
\toprule
Model & MRR (words) & Equiv.\ AUC (words) & MRR (digits) & Order probe (words) \\
\midrule
\vBase & \respm{\Ea}{stress_words.words_mrr} & \respm{\Ea}{stress_words.words_equivalence_auc} & \respm{\Ea}{retrieval.mrr} & \respm{\Ea}{ordering.words_probe_pairwise_order_acc} \\
\vParse & \respm{\Eb}{stress_words.words_mrr} & \respm{\Eb}{stress_words.words_equivalence_auc} & \respm{\Eb}{retrieval.mrr} & \respm{\Eb}{ordering.words_probe_pairwise_order_acc} \\
\vExact & \respm{\Ec}{stress_words.words_mrr} & \respm{\Ec}{stress_words.words_equivalence_auc} & \respm{\Ec}{retrieval.mrr} & \respm{\Ec}{ordering.words_probe_pairwise_order_acc} \\
\vExactParse & \respm{\Ee}{stress_words.words_mrr} & \respm{\Ee}{stress_words.words_equivalence_auc} & \respm{\Ee}{retrieval.mrr} & \respm{\Ee}{ordering.words_probe_pairwise_order_acc} \\
\vDual & \respm{\Ef}{stress_words.words_mrr} & \respm{\Ef}{stress_words.words_equivalence_auc} & \respm{\Ef}{retrieval.mrr} & \respm{\Ef}{ordering.words_probe_pairwise_order_acc} \\
\bottomrule
\end{tabular}}
\end{table}
\interp{the digit-to-words drop of \vParse and \vExactParse vs that of \vExact and \vDual. A large hybrid drop with a stable \vExact and \vDual is
evidence of learned numeracy; equal drops mean the gains were parser-driven.}

\subsection{Ablation}
\begin{table}[h]
\centering\small
\caption{Cumulative ablation (seed 42); the last row is \vExact.}
\label{tab:abl}
\begin{tabular}{@{}lcccccc@{}}
\toprule
Variant & Equiv.\ AUC & MRR & $\rho_{\log}$ & Order (probe) & Dim.\ discrim. & STS-B \\
\midrule
$\Lq_{\mathrm{sem}}$ only & \res{abl_A_semantic_only}{equivalence.equivalence_auc_hard} & \res{abl_A_semantic_only}{retrieval.mrr} & \res{abl_A_semantic_only}{distance.same_template_log_distance_spearman} & \res{abl_A_semantic_only}{ordering.probe_pairwise_order_acc} & \res{abl_A_semantic_only}{context.dimension_discrimination_acc} & \res{abl_A_semantic_only}{semantic.stsb_spearman} \\
$+\,\Lq_{\mathrm{eq}}$ & \res{abl_B_+equivalence}{equivalence.equivalence_auc_hard} & \res{abl_B_+equivalence}{retrieval.mrr} & \res{abl_B_+equivalence}{distance.same_template_log_distance_spearman} & \res{abl_B_+equivalence}{ordering.probe_pairwise_order_acc} & \res{abl_B_+equivalence}{context.dimension_discrimination_acc} & \res{abl_B_+equivalence}{semantic.stsb_spearman} \\
$+\,\Lq_{\mathrm{ord}}$ & \res{abl_C_+ordering}{equivalence.equivalence_auc_hard} & \res{abl_C_+ordering}{retrieval.mrr} & \res{abl_C_+ordering}{distance.same_template_log_distance_spearman} & \res{abl_C_+ordering}{ordering.probe_pairwise_order_acc} & \res{abl_C_+ordering}{context.dimension_discrimination_acc} & \res{abl_C_+ordering}{semantic.stsb_spearman} \\
$+\,\Lq_{\mathrm{dist}}$ & \res{abl_D_+distance}{equivalence.equivalence_auc_hard} & \res{abl_D_+distance}{retrieval.mrr} & \res{abl_D_+distance}{distance.same_template_log_distance_spearman} & \res{abl_D_+distance}{ordering.probe_pairwise_order_acc} & \res{abl_D_+distance}{context.dimension_discrimination_acc} & \res{abl_D_+distance}{semantic.stsb_spearman} \\
$+\,\Lq_{\mathrm{ctx}}$ (full) & \res{\Ec}{equivalence.equivalence_auc_hard} & \res{\Ec}{retrieval.mrr} & \res{\Ec}{distance.same_template_log_distance_spearman} & \res{\Ec}{ordering.probe_pairwise_order_acc} & \res{\Ec}{context.dimension_discrimination_acc} & \res{\Ec}{semantic.stsb_spearman} \\
\bottomrule
\end{tabular}
\end{table}
\interp{which step causes the largest change; single-seed steps smaller than the \vExact seed standard deviation are
inconclusive.}

\begin{figure}[h]
\centering
\begin{minipage}{0.48\linewidth}\centering\figorbox{03_equivalence_retrieval.png}{\linewidth}\end{minipage}\hfill
\begin{minipage}{0.48\linewidth}\centering\figorbox{06_pca_embedding_space.png}{\linewidth}\end{minipage}
\caption{Left: equivalence retrieval. Right: PCA of mass statements coloured by $\log_{10}$ magnitude, original vs
fine-tuned.}
\label{fig:retr}
\end{figure}

% =============================================================================================
\section{Discussion}
\label{sec:discussion}
% =============================================================================================
\paragraph{Higher cosine is not better numeracy.} A general-purpose encoder can assign high similarity to every pair
of template-sharing sentences. Equivalent pairs then score high, and so do pairs that differ by four orders of
magnitude or in dimension. Any evaluation that reports mean cosine on positives alone will mistake this compression
for understanding. The appropriate quantities are discriminative (AUC against hard negatives), graded (rank
correlation with numerical distance) and comparative (equivalent vs cross-dimension for the same anchor). Fine-tuning
can also \emph{lower} absolute cosines across the board while improving every one of these. That is the expected
signature of a space that has started to separate numbers, not a sign of degradation.

\paragraph{The verification paradox.} Verification is necessary, because an unverified small reasoner injects label
noise into exactly the relations it is supposed to teach. Verification also bounds the reasoner's value on
\emph{determinable} relations. After filtering, its equivalence, order and ratio labels coincide with what a few lines
of Python would produce. The reasoner can then contribute only (i) a \emph{selection effect}, since the pairs it gets
right may be systematically easier, and (ii) \emph{confidence weights}. The budget-matched control (\vMatched) is
therefore not a detail; without it, any gain of \vVerified over \vBase would be wrongly credited to distillation. If H4 holds,
the honest conclusion is that \emph{LLM reasoners are the wrong source of supervision for relations a parser can
compute}. Their value lies elsewhere: in plausibility (is 37\,\textdegree C a normal body temperature?), implicit units
(``a 6-footer''), approximate quantities (``about a dozen''), ranges and free-form phrasing. These are precisely the
signals Python cannot verify and a parser cannot extract.

\paragraph{Attribution: learned or parsed?} A parser hybrid will look near-perfect on any benchmark whose numbers are
written as digits with canonical units. That is a property of the benchmark, not evidence of learning. The words slice
breaks this: if the hybrid's advantage vanishes there while a fine-tuned or projected embedding keeps part of its gain,
the embedding has internalised some numeracy. Conversely, if all configurations collapse to baseline on spelled-out
numbers, the encoder learned digit-string regularities rather than quantities. That would be a negative result worth
reporting, and it points to number-aware tokenisation~\cite{thawani2021survey,sundararaman2020dice} as the bottleneck.

\paragraph{One space or two?} The tensions in \S\ref{sec:formulation} argue for decoupling. A learned numerical
subspace (\vDual) can make equivalent quantities collapse while the semantic space keeps contexts apart, and it does
this without a parser. Order, being antisymmetric, may ultimately require an asymmetric scoring
function~\cite{vendrov2016order} rather than any cosine. Whether one space is enough is thus an empirical question that
the \vExact vs \vDual comparison begins to answer.

\paragraph{Absolute vs logarithmic distance.} Logarithmic distance treats 1\,g vs 10\,g like 1\,kg vs 10\,kg, which suits
magnitude reasoning. Relative absolute distance treats 100\,g vs 110\,g as close and is closer to tolerance-based matching
in entity resolution. Neither is universally correct. A retrieval system should choose based on its downstream
decision, and a benchmark should report both, as we do.

\paragraph{Semantic cost.} Self-distillation on general text~\cite{reimers2020multilingual} is a cheap guard against
forgetting. Our numerical objectives, however, deliberately reshape similarities between near-identical sentences, so
some movement on STS-B is expected. The relevant question is the size of the trade-off, not whether one exists.

% =============================================================================================
\section{Limitations and Threats to Validity}
% =============================================================================================
\textbf{Synthetic, template-based data.} Real text has implicit units, ranges, approximations, multiple quantities per
sentence and domain jargon. Outside category G, values are deliberately decoupled from plausibility (for example a
12\,m\textsuperscript{3} syrup dose) to isolate numerical relations. \textbf{Templates.} Templates are shared across
splits by design, except for seven held-out templates; the words slice tests format shift, not new vocabulary.
\textbf{Reasoner scale.} Conclusions about a 1.5B (or 0.5B) reasoner need not transfer to larger models.
\textbf{Hyperparameters} ($\lambda$, margins, $\sigma$) were set a priori, not tuned per configuration.
\textbf{Probes} measure linear decodability, not downstream use~\cite{hewitt2019control}. \textbf{Statistics.} Three
seeds for the main configurations and one for the ablation; bootstrap intervals cover test-item sampling only.

% =============================================================================================
\section{Conclusion and Research Agenda}
% =============================================================================================
Transferring numerical reasoning into sentence embeddings is attractive and easy to evaluate badly. We have argued that
three controls are necessary to make claims about it: discriminative rather than absolute similarity, a parser-free
test slice, and a budget-matched reasoner-free baseline. We have built a benchmark and protocol that implement them, with
pre-declared hypotheses. \interp{one-paragraph summary of which hypotheses were supported.}

We see four directions as most promising: (i) using LLM reasoners for \emph{non-determinable} numerical judgements
(plausibility, approximation, implicit units, ranges), where distillation can genuinely add information; (ii)
asymmetric or box-style heads for order; (iii) number-aware tokenisation inside the encoder; and (iv) evaluation on
real numeric retrieval and entity-resolution tasks, such as product specifications, clinical notes and scientific
tables.

\paragraph{Reproducibility.} Code, configuration and the single entry point (\texttt{main.py}) are available at
\url{https://github.com/TariqMahmood93/Numerical_Intelligence_Transfor}.

\bibliographystyle{plainnat}
\bibliography{refs}

\appendix
% =============================================================================================
\section{Hyperparameters}
\begin{table}[h]
\centering\small
\resizebox{\linewidth}{!}{%
\begin{tabular}{@{}ll@{\hspace{2em}}ll@{}}
\toprule
Encoder & all-MiniLM-L6-v2 & Reasoner & Qwen2.5-1.5B-Instruct ($\ge$8\,GB GPU) \\
Optimiser & AdamW, lr $2\times10^{-5}$, wd 0.01 & Schedule & 10\,\% warm-up, linear decay \\
Epochs $\times$ steps & $4\times300$ & Batch & 32 pairs (16 triplets) per objective \\
$\lambda_{1..5}$ & 1.0, 0.5, 1.0, 0.5, 1.0 & $\tau$ (InfoNCE) & 0.05 \\
$m_r$, $m_t$ & 0.10, 0.05 & $m_{\mathrm{cross}}, m_{\mathrm{ctx}}, m_{\mathrm{plaus}}$ & 0.30, 0.85, 0.70 \\
$c_{\min}, c_{\max}$ & 0.2, 0.9 & $\sigma_{\log}$, $\sigma_{\mathrm{abs}}$ & 1.0, 0.25 \\
Projection & $384\to128$ & Reasoner budget & 4{,}000 pairs (GPU) \\
Seeds & 42, 43, 44 & Model selection & unweighted validation loss \\
\bottomrule
\end{tabular}}
\end{table}

\section{Reasoner Prompts}
\label{app:prompts}
System: \emph{``You are a careful assistant for physical quantities. Convert units exactly before comparing.
Temperatures must be compared on an absolute scale. Answer with only the requested option.''}
\begin{itemize}[leftmargin=*,itemsep=1pt]
\item \textbf{Dimension:} ``Do A and B measure the same physical dimension (\ldots)? Answer Yes or No.''
\item \textbf{Equivalence:} ``After converting to a common unit, do A and B state exactly the same quantity? Answer Yes or No.''
\item \textbf{Order:} ``After converting to a common unit, which quantity is larger? Answer A, B, or Equal.''
\item \textbf{Ratio:} ``Convert both quantities to the same unit (use kelvin for temperatures) and compute the ratio A
divided by B. Reply with only the number.''
\item \textbf{Plausibility:} ``Is the numerical value physically plausible given the description in the statement?
Answer Yes or No.''
\end{itemize}

\end{document}
